Signed-distance-field local geometry feedforward offers a low-intrusion way to strengthen complementarity-based rigid-body contact. Nonsmooth contact dynamics based on velocity-level time-stepping schemes and linear complementarity problems (LCPs) remains one of the core paradigms for rigid-body collision, intermittent contact, and contact-state transitions in robotics simulation, computer graphics, and engineering multibody dynamics \cite{flores2022contact,wang2021nonsmooth,stewart1996implicit,anitescu1997formulating}. Within this class of methods, complementarity constraints provide a natural way to impose non-penetration, impact response, and contact-state switching, which is why they remain particularly attractive for intermittent rigid contact. Practical contact quality, however, is determined not only by the algebraic form of the constraints but also by how local geometric information enters the model.

Conventional contact detection usually relies on polygon meshes and boundary-distance queries, with GJK and its descendants serving as representative examples \cite{gilbert1988fast,wang2021collision}. In recent years, the emergence of sparse voxel storage, GPU-friendly data structures, and SDF-based collision handling has made the use of signed distance fields (SDFs) as implicit surrogates for rigid-body surfaces increasingly attractive \cite{museth2021nanovdb,macklin2020sdfcollision,liu2024sdfsdf,lai2024unifying,lopez2024convexsdf}. However, within a complementarity-based rigid-contact framework, if local geometry enters only through the distance value $\Phi$ and the first-order gradient $\nabla \Phi$, the contact normal, gap evolution, and constraint response can become highly sensitive to local geometric error during contact onset, sustained sliding, and local state transitions. In geometrically demanding scenarios such as cam tips, roller transition zones, and gear boundaries, this sensitivity may manifest itself as velocity spikes, peak errors, and degraded tracking behavior.

In this setting, the focus of the present work is not the redesign of the underlying complementarity solver, but the input-layer modeling problem that arises when SDF geometry is coupled to an existing solver chain. In many engineering cases, the limiting factor is whether local geometric information enters the contact model in a sufficiently stable and accurate manner, especially in the prediction of gap evolution, the description of local normal manifolds, and the control of geometric noise during contact activation. From this perspective, a low-intrusion enhancement of the geometry-to-solver interface within an established LCP framework is more aligned with the problem structure addressed here than the construction of yet another rigid-contact solver.

It is important to stress that the present work does not attempt to extend complementarity-based methods into regimes that are better handled by penalty-based formulations, such as soft contact or complex distributed contact. Instead, the focus is a more basic modeling question within the valid scope of complementarity-based rigid contact: \textbf{without changing the existing hard-contact complementarity solution framework and without switching to penalty-based soft-contact modeling, how should local contact geometry be represented so that more accurate contact responses can be obtained during contact onset, sustained contact, and state transitions without sacrificing computational efficiency?} In other words, the question is not whether one enhanced geometric proxy can universally outperform traditional approaches, but rather how local geometric information should enter a complementarity-based rigid-contact model in a way that matches the contact state and converts geometric detail into simultaneous gains in numerical accuracy and computational efficiency under clearly defined applicability conditions.

To address this question, this paper develops a \textbf{local-geometry-aware complementarity-based rigid-contact modeling strategy}. Within a unified Moreau--Jean time-stepping framework, local geometric information from the SDF is used to construct local representations in the neighborhood of contact, and the representation is configured according to the contact state. For cases with continuous local contact evolution and relatively smooth sliding, a \texttt{SlidingPatch} representation with a persistent manifold is adopted, and Hessian-based curvature feedforward is injected on the corresponding local manifold. For cases characterized by contact onset, compact local contact domains, or sensitivity to state transitions, a more compact and conservative \texttt{CompactDiscrete} representation is used to preserve numerical robustness. In the current implementation, this selection is performed through case- or application-level preset configuration rather than runtime switching by a unified classifier. Accordingly, the main purpose of this work is neither to replace the underlying solver nor to propose an automatic regime-identification algorithm, but to analyze whether regime-matched local geometric representations can improve numerical accuracy without sacrificing computational efficiency inside an existing complementarity-based rigid-contact framework.

The main contributions of this work are as follows:
\begin{itemize}
    \item \textbf{A local-geometry-aware complementarity model for rigid contact.} Local SDF-based geometric information is introduced at the contact-modeling level without replacing the mature global solver backbone. This places the improvement explicitly at the geometry-to-solver interface, enriches the geometric description near the contact region, and supports accurate contact simulation without sacrificing computational efficiency within the intended application regimes.
    \item \textbf{A contact-state-oriented configuration strategy for local representations.} For continuous sliding, a continuous \texttt{SlidingPatch} representation is combined with local curvature feedforward; for contact onset, compact contact domains, and state-transition-sensitive cases, a compact discrete \texttt{CompactDiscrete} representation is adopted, so that the use of local geometry is grounded in a representation matched to the contact regime.
    \item \textbf{A systematic validation of the applicability of local-geometry-aware modeling using existing benchmark cases.} Through head-on impact, an eccentric disk--roller analytical benchmark, the Onset Stress Benchmark, an industrial cam, and a Twin Gear case, the proposed modeling strategy is evaluated across different contact scenarios. The results show that, under regime-matched settings, it improves local geometric description and global numerical response in continuous sliding and contact-onset problems while maintaining computational efficiency; in the presence of local singular boundaries or nonsmooth meshing/contact transitions, its effectiveness depends more strongly on the match between local representation and contact state.
\end{itemize}

It should also be emphasized that the experiments in this paper are not intended to prove that one single regime is universally optimal across all benchmarks. Instead, each benchmark is regarded as a representative scenario associated with a particular contact state, and the proposed modeling strategy is assessed under a preconfigured regime matched to its local geometric characteristics so as to examine whether stable, interpretable, and practically meaningful gains in both efficiency and accuracy can be obtained. The ultimate outcome is therefore a modeling principle for regime-matched contact simulation rather than a runtime algorithm that automatically selects the best local representation.
